\scanpage{179}
Ceci dit, pour définir \inserted{un} $f$ \deleted{$\unsure{f^\circ}$} correspondant à un $f^\circ$ déjà donné, on peut se le donner arbitrairement \inserted{\unsure{(soit $g$)}} sur $\prod F_i$ (comme applic\unsure{ation} multilin.).

\deleted{On \unsure{aura donc}}
Faisons le calcul pour deux facteurs ---
On pose (avec $x_i\in F_i$, $\alpha_i\in E_i^\circ$)
\[
\begin{aligned}
f(\underbrace{x_1+\alpha_1}_{X_1},\underbrace{x_2+\alpha_2}_{X_2})
 &=g(x_1,x_2)+\deleted{\unsure{f(x_1,\alpha_2)}}f^\circ(\alpha_1,b_2(x_2))\\
 &\quad+f^\circ(b_1(x_1),\alpha_2)
 +\underbrace{\lambda f^\circ(\alpha_1,\alpha_2)}_{\substack{\Vert\\f^\circ(u_1(\alpha_1),\alpha_2)\\=f^\circ(\alpha_1,u_2(\alpha_2))}}
\end{aligned}
\]
Il suffit de vérifier que ça satisfait aux deux conditions. P. ex si on remplace $\alpha_1$ \deleted{\unsure{par}} $\alpha_1+\alpha'_1$, on obtient
\[
\begin{aligned}
f(X_1+\alpha'_1,X_2)&=f(X_1,X_2)+f^\circ(\alpha'_1,\uline{b_2(x_2)})\\
&\qquad+\lambda f^\circ(\alpha'_1,\alpha_2)
\end{aligned}
\]
et on \unsure{veut}
\[
=f(\unsure{x_1|X_1},X_2)+f^\circ\bigl(\alpha'_1,\underbrace{u_2(X_2)}_{b_2(x_2)+\lambda\alpha_2}\bigr)\qquad\text{OK.}
\]
Prenant pour $B$ le produit tensoriel $E_1^\circ\otimes E_2^\circ$, pour $f^\circ$ l'accouplement can., et $g=0$, on trouve
\margin{178}
